% Hahn分解的非唯一性：eta=lambda|_[0,1]-lambda|_[2,3]，块高表示长度测度前的系数(+1,-1)；两种Hahn分解 P1=(-inf,3/2)、P2=(-inf,2)，差集[3/2,2)上eta的所有子集取值为零。横轴1单位=1.6cm。
\begin{tikzpicture}[
  lab/.style={font=\figurefont\fontsize{8}{10}\selectfont,text=FigureInk},
  note/.style={font=\FigureNoteFont\fontsize{7.5}{9}\selectfont,text=FigureMuted,align=center},
  ax/.style={draw=FigureStroke,line width=.75pt,-{Stealth[length=2.2mm]}},
  posbar/.style={draw=FigureBlue,fill=FigureBlueFill,line width=.7pt},
  negbar/.style={draw=FigureRed,fill=FigureRedFill,line width=.7pt}]
\def\u{1.6}
\fill[FigureYellowFill] ({1.5*\u},-.75) rectangle ({2*\u},.75);
\draw[posbar] (0,0) rectangle (\u,.6);
\draw[negbar] ({2*\u},0) rectangle ({3*\u},-.6);
\draw[ax] (-1,0)--({3*\u+.9},0);
\foreach \x in {0,1,2,3}{\draw[FigureStroke,line width=.6pt] ({\x*\u},0)--({\x*\u},-.08);}
\node[lab,below] at (0,-.06) {$0$}; \node[lab,below] at (\u,-.06) {$1$};
\node[lab,above] at ({2*\u},.02) {$2$}; \node[lab,above] at ({3*\u},.02) {$3$};
\node[lab] at ({\u/2},.95) {$+\lambda|_{[0,1]}$};
\node[lab] at ({2.5*\u},-.95) {$-\lambda|_{[2,3]}$};
% two Hahn choices
\draw[FigureBlue,line width=.9pt,-{Stealth[length=1.8mm]}] ({1.5*\u},1.55)--(-1,1.55);
\draw[FigureBlue,line width=.9pt] ({1.5*\u},1.45)--({1.5*\u},1.65);
\node[lab,anchor=south] at ({0.25*\u},1.6) {$P_1=(-\infty,3/2)$};
\draw[FigureGreen,line width=.9pt,-{Stealth[length=1.8mm]}] ({2*\u},-1.55)--(-1,-1.55);
\draw[FigureGreen,line width=.9pt] ({2*\u},-1.45)--({2*\u},-1.65);
\node[lab,anchor=north] at ({0.5*\u},-1.6) {$P_2=(-\infty,2)$};
\node[note,anchor=west] at ({2.05*\u},1.25) {$[3/2,2)$：$\eta$在其每个子集上为零};
\end{tikzpicture}
